\heading{第12章-跋}
\begin{question}
纠缠态（Entangled states）。纠缠态的一个经典例子是自旋单态（式 (12.1)）：两个粒子体系的状态不能写成两个单粒子状态波函数的乘积。考虑两能级体系 $|\phi_a\rangle$、$|\phi_b\rangle$，且 $\langle\phi_i|\phi_j\rangle=\delta_{ij}$。证明：对于任意单粒子状态 $|\psi_r\rangle$ 和 $|\psi_s\rangle$，两粒子体系状态
\[
\alpha|\phi_a(1)\rangle|\phi_b(2)\rangle+\beta|\phi_b(1)\rangle|\phi_a(2)\rangle\qquad(\alpha\ne0,\ \beta\ne0)
\]
都不能表示为乘积形式 $|\psi_r(1)\rangle|\psi_s(2)\rangle$。提示：把 $|\psi_r\rangle$ 和 $|\psi_s\rangle$ 分别写成 $|\phi_a\rangle$ 和 $|\phi_b\rangle$ 的线性组合。

\par\smallskip
\noindent{\kaishu 为避免查阅正文，本题中引用的公式为：式 (12.1)：两个自旋 $1/2$ 的单态：$|0,0\rangle=(|\uparrow\downarrow\rangle-|\downarrow\uparrow\rangle)/\sqrt2$，不能分解为两个单粒子态的乘积。}

\end{question}
\begin{solution}
假设该两粒子态可以写成乘积形式
\[
|\Psi\rangle=|\psi_r(1)\rangle|\psi_s(2)\rangle,
\]
其中 $|\psi_r\rangle$ 和 $|\psi_s\rangle$ 是单粒子态。在正交归一基 $\{|\phi_a\rangle,|\phi_b\rangle\}$ 中展开这两个单粒子态：
\[
|\psi_r\rangle=c_a|\phi_a\rangle+c_b|\phi_b\rangle,\qquad
|\psi_s\rangle=d_a|\phi_a\rangle+d_b|\phi_b\rangle,
\]
且满足归一化条件
\[
|c_a|^2+|c_b|^2=1,\qquad |d_a|^2+|d_b|^2=1.
\]
将乘积形式展开：
\[
\begin{aligned}
|\psi_r(1)\rangle|\psi_s(2)\rangle
&=\bigl(c_a|\phi_a(1)\rangle+c_b|\phi_b(1)\rangle\bigr)\bigl(d_a|\phi_a(2)\rangle+d_b|\phi_b(2)\rangle\bigr)\\[2mm]
&=c_a d_a\,|\phi_a(1)\rangle|\phi_a(2)\rangle
+ c_a d_b\,|\phi_a(1)\rangle|\phi_b(2)\rangle\\
&\quad + c_b d_a\,|\phi_b(1)\rangle|\phi_a(2)\rangle
+ c_b d_b\,|\phi_b(1)\rangle|\phi_b(2)\rangle.
\end{aligned}
\]
为简便记 $|aa\rangle\equiv|\phi_a(1)\rangle|\phi_a(2)\rangle$，$|ab\rangle\equiv|\phi_a(1)\rangle|\phi_b(2)\rangle$，等等，则
\[
|\psi_r(1)\rangle|\psi_s(2)\rangle
=c_a d_a|aa\rangle + c_a d_b|ab\rangle + c_b d_a|ba\rangle + c_b d_b|bb\rangle.
\]

题给态为
\[
|\Psi\rangle=\alpha|ab\rangle+\beta|ba\rangle,
\]
其中 $\alpha\neq0$, $\beta\neq0$。如果 $|\Psi\rangle$ 确实可以写成乘积形式，则比较以上两式对应基矢的系数，得到方程组
\[
\begin{cases}
c_a d_a = 0, & \text{（$|aa\rangle$ 项的系数）}\\[2pt]
c_a d_b = \alpha, & \text{（$|ab\rangle$ 项的系数）}\\[2pt]
c_b d_a = \beta, & \text{（$|ba\rangle$ 项的系数）}\\[2pt]
c_b d_b = 0. & \text{（$|bb\rangle$ 项的系数）}
\end{cases}
\]

由第二式 $c_a d_b = \alpha$，因 $\alpha\neq0$，必有 $c_a\neq0$ 且 $d_b\neq0$。
由第三式 $c_b d_a = \beta$，因 $\beta\neq0$，必有 $c_b\neq0$ 且 $d_a\neq0$。
然而第一式 $c_a d_a = 0$ 要求 $c_a$ 和 $d_a$ 至少有一个为零，这与 $c_a\neq0$, $d_a\neq0$ 矛盾。
同时，第四式 $c_b d_b = 0$ 要求 $c_b$ 和 $d_b$ 至少有一个为零，这与 $c_b\neq0$, $d_b\neq0$ 矛盾。

因此该方程组无解，假设不成立。故对于任意非零的 $\alpha,\beta$，态
\[
\boxed{|\Psi\rangle=\alpha|\phi_a(1)\rangle|\phi_b(2)\rangle+\beta|\phi_b(1)\rangle|\phi_a(2)\rangle}
\]
都不能写成单粒子态的乘积形式，是纠缠态。
\end{solution}

\begin{question}
爱因斯坦盒子（Einstein's Boxes）。想象一个粒子被限制在盒子里（可看作一维无限方势阱），处于基态；当引入一个不可穿透隔板时，它将盒子分成 $B_1$ 和 $B_2$ 两半。随后将两部分移开相距很远，并对 $B_1$ 进行测量查看粒子是否在其中。若答案为肯定，立刻知道粒子不在远处的 $B_2$ 中。
\begin{enumerate}
\item 爱因斯坦会怎么说明这一点？
\item 哥本哈根学派如何解释？在对 $B_1$ 测量之后，$B_2$ 中的波函数是什么？
\end{enumerate}
\end{question}
\begin{solution}
(a) 爱因斯坦从"物理实在性"的立场出发认为：粒子在测量之前已经客观地位于某一个半盒中。对 $B_1$ 的测量只是揭示这个既有事实；远处 $B_2$ 的物理状态并没有被瞬时的超距作用改变。爱因斯坦将此看作波函数描述不完备的证据——如果波函数给出了粒子在两个半盒中的叠加，而实际粒子只在一个半盒中，那么波函数就没有完全反映物理实在。

(b) 哥本哈根解释中，在隔板插入之后、测量之前，粒子态是两个空间局域波包的等权重叠加：
\[
|\Psi\rangle={1\over\sqrt2}\bigl(|B_1\rangle+|B_2\rangle\bigr),
\]
其中 $|B_1\rangle$ 和 $|B_2\rangle$ 分别表示粒子完全局域在左半盒和右半盒的波函数，二者空间区域不重叠。这里忽略可能的相对相位。

对 $B_1$ 进行测量后，根据投影公理：
\begin{itemize}
\item 若测量结果为"有粒子"，态立即坍缩为 $|B_1\rangle$，此时 $B_2$ 中的波函数为零（粒子不在 $B_2$ 中），即
\[
|\Psi\rangle_{\text{后}} = |B_1\rangle,\qquad \Psi_{B_2}(x)=0.
\]
\item 若测量结果为"无粒子"，态立即坍缩为 $|B_2\rangle$，此时 $B_2$ 中的波函数恢复为归一化的盒中基态波函数：
\[
|\Psi\rangle_{\text{后}} = |B_2\rangle,\qquad \Psi_{B_2}(x) = \sqrt{\frac{2}{L}}\sin\frac{\pi x}{L},
\]
其中 $L$ 为半盒宽度。
\end{itemize}

因此，在对 $B_1$ 测量之后，$B_2$ 中的波函数要么为零（肯定结果），要么为完整的归一化单盒基态（否定结果）。这一"瞬时坍缩"正是哥本哈根解释中波函数更新规则的直接结果，也是 EPR 悖论的核心争议所在。
\end{solution}

\begin{question}
（局域）确定性（"隐变量"）理论的一个例子是经典力学。设用宏观物体（如棒球）代替电子和质子进行贝尔实验，并令隐变量 $\lambda=(\theta,\phi)$ 表示角动量方向。探测器只记录沿 $\mathbf a$ 和 $\mathbf b$ 方向的分量符号：
\[
A=\operatorname{sign}(\mathbf a\cdot\mathbf S_a),\qquad B=\operatorname{sign}(\mathbf b\cdot\mathbf S_b).
\]
\begin{enumerate}
\item 取图 12.5 中的坐标轴，使 $\mathbf a$ 与 $\mathbf b$ 位于 $x$-$y$ 平面内，$\mathbf a$ 沿 $x$ 轴。证明
\[
A(\mathbf a,\lambda)B(\mathbf b,\lambda)=-\operatorname{sign}[\cos\phi\cos(\phi-\eta)],
\]
其中 $\eta$ 为 $\mathbf a$ 和 $\mathbf b$ 之间的夹角。
\item 假设棒球发射方向的可能性均匀，计算 $P(\mathbf a,\mathbf b)$。
\item 在 $\eta=-\pi$ 到 $+\pi$ 的范围内作图 $P(\mathbf a,\mathbf b)$，并在同一图上画出量子公式；取 $\eta$ 的什么值时，隐变量理论与量子力学结果一致？
\item 验证你的结果满足式 (12.12) 所表示的贝尔不等式。
\end{enumerate}

\par\smallskip
\noindent{\kaishu 为避免查阅正文，本题中引用的公式为：式 (12.12)：Bell 型不等式的相关函数形式：$|P(\mathbf a,\mathbf b)-P(\mathbf a,\mathbf c)|\le 1+P(\mathbf b,\mathbf c)$。}

\end{question}
\begin{solution}
(a) 取坐标系使 $\mathbf a$ 沿 $x$ 轴方向，即 $\mathbf a = \hat x$。设 $\mathbf b$ 在 $x$-$y$ 平面内与 $\mathbf a$ 成夹角 $\eta$，于是
\[
\mathbf a = \hat x,\qquad \mathbf b = \cos\eta\,\hat x + \sin\eta\,\hat y.
\]

对于经典棒球，两粒子的角动量方向相反（总角动量为零），设粒子1的角动量方向为单位矢量 $\hat S$，其球坐标为 $(\theta,\phi)$，则
\[
\hat S = \sin\theta\cos\phi\,\hat x + \sin\theta\sin\phi\,\hat y + \cos\theta\,\hat z.
\]
粒子2的角动量方向为 $-\hat S$。

探测器 A 沿 $\mathbf a$ 方向测量粒子1的角动量分量符号：
\[
A(\mathbf a,\lambda) = \operatorname{sign}(\mathbf a\cdot\hat S) = \operatorname{sign}(\sin\theta\cos\phi).
\]
探测器 B 沿 $\mathbf b$ 方向测量粒子2的角动量分量符号：
\[
B(\mathbf b,\lambda) = \operatorname{sign}(\mathbf b\cdot(-\hat S)) = -\operatorname{sign}(\mathbf b\cdot\hat S).
\]
计算内积：
\[
\mathbf b\cdot\hat S = \cos\eta\sin\theta\cos\phi + \sin\eta\sin\theta\sin\phi = \sin\theta\cos(\phi-\eta).
\]
因此
\[
B(\mathbf b,\lambda) = -\operatorname{sign}[\sin\theta\cos(\phi-\eta)].
\]
对于 $\sin\theta > 0$（极角 $\theta$ 不在极点），符号函数中 $\sin\theta$ 为正因子可不计入，从而
\[
A(\mathbf a,\lambda) = \operatorname{sign}(\cos\phi),\qquad
B(\mathbf b,\lambda) = -\operatorname{sign}[\cos(\phi-\eta)].
\]
两测量结果的乘积为
\[
\boxed{A(\mathbf a,\lambda)B(\mathbf b,\lambda) = -\operatorname{sign}[\cos\phi\cos(\phi-\eta)].}
\]

(b) 假设隐变量 $\lambda$ 均匀分布，即 $\phi$ 在 $[0,2\pi)$ 上均匀分布，$\theta$ 的分布不影响符号函数结果（只要 $\sin\theta>0$）。相关函数 $P(\mathbf a,\mathbf b)$ 定义为 $A$ 和 $B$ 的系综平均：
\[
P(\mathbf a,\mathbf b) = \langle AB\rangle = \frac{1}{2\pi}\int_0^{2\pi} A(\mathbf a,\phi)B(\mathbf b,\phi)\,\dd\phi.
\]

根据 (a) 的结果，$AB = -1$ 当 $\cos\phi$ 与 $\cos(\phi-\eta)$ 同号；$AB = +1$ 当两者异号。我们需要计算 $\cos\phi$ 和 $\cos(\phi-\eta)$ 同号的 $\phi$ 区间总长度。

在一个周期 $2\pi$ 内，$\cos\phi$ 的符号在 $\phi = \pi/2$ 和 $3\pi/2$ 处改变；$\cos(\phi-\eta)$ 的符号在 $\phi = \pi/2+\eta$ 和 $3\pi/2+\eta$ 处改变。考虑 $0\le |\eta|\le\pi$ 的情形：

两个余弦同号的 $\phi$ 区间总长度为 $2(\pi-|\eta|)$，异号的区间总长度为 $2|\eta|$。因此
\[
\langle AB\rangle = \frac{1}{2\pi}\bigl[(-1)\cdot 2(\pi-|\eta|) + (+1)\cdot 2|\eta|\bigr]
= -\frac{2\pi-4|\eta|}{2\pi} = -1 + \frac{2|\eta|}{\pi}.
\]
故
\[
\boxed{P(\mathbf a,\mathbf b) = -1 + \frac{2|\eta|}{\pi},\qquad 0\le|\eta|\le\pi.}
\]

(c) 量子力学中，自旋 $1/2$ 单态的相关函数为
\[
\boxed{P_Q(\mathbf a,\mathbf b) = -\cos\eta.}
\]
比较经典隐变量结果 $P_{\text{cl}} = -1+2|\eta|/\pi$ 与量子结果 $P_Q = -\cos\eta$：

令 $P_{\text{cl}}(\eta)=P_Q(\eta)$，即
\[
-1+\frac{2|\eta|}{\pi} = -\cos\eta.
\]
求解可得交点位于：
\begin{itemize}
\item $\eta=0$：$P_{\text{cl}}(0) = -1$，$P_Q(0) = -1$。
\item $\eta=\pm\pi/2$：$P_{\text{cl}}(\pm\pi/2) = -1+1 = 0$，$P_Q(\pm\pi/2) = 0$。
\item $\eta=\pm\pi$：$P_{\text{cl}}(\pm\pi) = -1+2 = 1$，$P_Q(\pm\pi) = 1$。
\end{itemize}
此外两个函数在其它 $\eta$ 处不一致。在 $0<\eta<\pi/2$ 区域，$P_Q < P_{\text{cl}}$（量子关联更强）；在 $\pi/2<\eta<\pi$ 区域，$P_Q > P_{\text{cl}}$。

(d) 贝尔不等式（式 (12.12)）为
\[
|P(\mathbf a,\mathbf b)-P(\mathbf a,\mathbf c)| \le 1+P(\mathbf b,\mathbf c).
\]
由于该经典棒球模型中，每个粒子对都具有确定的局域函数 $A(\mathbf a,\lambda)$ 和 $B(\mathbf b,\lambda)$，取值仅为 $\pm1$，对所有 $\lambda$ 有
\[
|A(\mathbf a,\lambda)B(\mathbf b,\lambda)-A(\mathbf a,\lambda)B(\mathbf c,\lambda)| = |A(\mathbf a,\lambda)B(\mathbf b,\lambda)[1-B(\mathbf b,\lambda)B(\mathbf c,\lambda)]|
\]
因为 $A(\mathbf a,\lambda)B(\mathbf b,\lambda)=\pm1$ 以及 $B(\mathbf b,\lambda)B(\mathbf c,\lambda)=\pm1$，上式绝对值即为
\[
|1-B(\mathbf b,\lambda)B(\mathbf c,\lambda)| \le 1+B(\mathbf b,\lambda)B(\mathbf c,\lambda).
\]
对所有 $\lambda$ 取系综平均（注意 $B^2=1$）：
\[
|P(\mathbf a,\mathbf b)-P(\mathbf a,\mathbf c)| \le 1+P(\mathbf b,\mathbf c).
\]
因此这一经典隐变量模型自然满足贝尔不等式。代入 (b) 的结果：
\[
P(\mathbf a,\mathbf b) = -1+\frac{2|\eta_{ab}|}{\pi},
\]
取 $\eta_{ab}=\eta$，$\eta_{ac}=2\eta$，$\eta_{bc}=\eta$（例如 $\mathbf a$、$\mathbf b$、$\mathbf c$ 依次排列），容易验证不等式成立。
\end{solution}

\begin{question}
\begin{enumerate}
\item 证明式 (12.17)--式 (12.20)。
\item 证明密度算符随时间演化由下式确定：
\[
i\hbar\frac{\dd\hat\rho}{\dd t}=[\hat H,\hat\rho].
\]
这是由 $\hat\rho$ 表示的薛定谔方程。
\end{enumerate}

\par\smallskip
\noindent{\kaishu 为避免查阅正文，本题中引用的公式为：式 (12.17--12.20)：密度算符：$\rho=\sum_i p_i|\psi_i\rangle\langle\psi_i|$，$\rho^\dagger=\rho$，$\operatorname{Tr}\rho=1$，$\langle Q\rangle=\operatorname{Tr}(\rho Q)$，纯态为 $\rho^2=\rho$。}

\end{question}
\begin{solution}
(a) 首先给出密度算符的定义。设系统以概率 $p_i$ 处于态 $|\psi_i\rangle$（$p_i\ge0$，$\sum_i p_i=1$，$\{|\psi_i\rangle\}$ 不一定正交），则密度算符为
\[
\rho = \sum_i p_i |\psi_i\rangle\langle\psi_i|.
\]
这就是式 (12.17)。

\textbf{式 (12.18)：厄米性。}
取厄米共轭：
\[
\rho^\dagger = \Bigl(\sum_i p_i |\psi_i\rangle\langle\psi_i|\Bigr)^\dagger
= \sum_i p_i (|\psi_i\rangle\langle\psi_i|)^\dagger
= \sum_i p_i |\psi_i\rangle\langle\psi_i| = \rho.
\]
因此 $\rho^\dagger = \rho$，即密度算符是厄米算符。

\textbf{式 (12.19)：迹为 1。}
在任意正交归一基 $\{|n\rangle\}$ 中计算迹：
\[
\operatorname{Tr}\rho = \sum_n \langle n|\rho|n\rangle
= \sum_n \sum_i p_i \langle n|\psi_i\rangle\langle\psi_i|n\rangle
= \sum_i p_i \sum_n \langle\psi_i|n\rangle\langle n|\psi_i\rangle
= \sum_i p_i \langle\psi_i|\psi_i\rangle = \sum_i p_i = 1.
\]
这里利用了 $\sum_n |n\rangle\langle n| = I$ 的完备性关系以及 $\langle\psi_i|\psi_i\rangle=1$。

\textbf{式 (12.20)：期望值。}
任意可观测量 $Q$ 的期望值为概率加权平均：
\[
\langle Q\rangle = \sum_i p_i \langle\psi_i|Q|\psi_i\rangle.
\]
利用迹的循环性质 $\operatorname{Tr}(AB)=\operatorname{Tr}(BA)$，可将其改写为
\[
\langle Q\rangle = \sum_i p_i \langle\psi_i|Q|\psi_i\rangle
= \sum_i p_i \operatorname{Tr}\bigl(Q|\psi_i\rangle\langle\psi_i|\bigr)
= \operatorname{Tr}\Bigl(Q\sum_i p_i|\psi_i\rangle\langle\psi_i|\Bigr)
= \operatorname{Tr}(\rho Q).
\]

\textbf{纯态条件 $\rho^2=\rho$。}
对于纯态 $\rho=|\psi\rangle\langle\psi|$，
\[
\rho^2 = |\psi\rangle\langle\psi|\psi\rangle\langle\psi| = |\psi\rangle\langle\psi| = \rho,
\]
因为 $\langle\psi|\psi\rangle=1$。反之，若 $\rho^2=\rho$ 且 $\operatorname{Tr}\rho=1$，则 $\rho$ 是投影算符且秩为 1，代表纯态。

\textbf{基变换下的变换规律。}
设基矢变换 $|n'\rangle = U|n\rangle$，其中 $U$ 是幺正算符。密度矩阵的矩阵元变换为
\[
\rho'_{mn'} = \langle m'|\rho|n'\rangle = \langle m|U^\dagger\rho U|n\rangle = (U^\dagger\rho U)_{mn},
\]
即 $\rho' = U^\dagger\rho U$，等价于 $\rho' = U\rho U^\dagger$（约定不同可相差一个转置）。

(b) 首先考虑纯态情况。设 $|\psi(t)\rangle$ 满足薛定谔方程
\[
i\hbar\frac{\dd}{\dd t}|\psi(t)\rangle = H|\psi(t)\rangle,
\]
其厄米共轭为
\[
-i\hbar\frac{\dd}{\dd t}\langle\psi(t)| = \langle\psi(t)|H.
\]
纯态密度算符 $\rho(t)=|\psi(t)\rangle\langle\psi(t)|$ 对时间求导：
\[
\begin{aligned}
\frac{\dd\rho}{\dd t}
&= \Bigl(\frac{\dd}{\dd t}|\psi\rangle\Bigr)\langle\psi| + |\psi\rangle\Bigl(\frac{\dd}{\dd t}\langle\psi|\Bigr)\\[2mm]
&= \Bigl(\frac{1}{i\hbar}H|\psi\rangle\Bigr)\langle\psi| + |\psi\rangle\Bigl(-\frac{1}{i\hbar}\langle\psi|H\Bigr)\\[2mm]
&= \frac{1}{i\hbar} H|\psi\rangle\langle\psi| - \frac{1}{i\hbar} |\psi\rangle\langle\psi|H \\[2mm]
&= \frac{1}{i\hbar}(H\rho - \rho H) = \frac{1}{i\hbar}[H,\rho].
\end{aligned}
\]
两边乘以 $i\hbar$，得
\[
\boxed{i\hbar\frac{\dd\rho}{\dd t} = [H,\rho].}
\]

对于混合态 $\rho = \sum_i p_i |\psi_i\rangle\langle\psi_i|$，其中 $p_i$ 为与时间无关的经典概率（系综混合，而非量子叠加），每个纯态投影 $|\psi_i(t)\rangle\langle\psi_i(t)|$ 独立地满足上述方程：
\[
\frac{\dd}{\dd t}|\psi_i\rangle\langle\psi_i| = \frac{1}{i\hbar}[H,|\psi_i\rangle\langle\psi_i|].
\]
取概率加权和：
\[
\frac{\dd\rho}{\dd t} = \sum_i p_i\frac{\dd}{\dd t}|\psi_i\rangle\langle\psi_i|
= \frac{1}{i\hbar}\sum_i p_i [H,|\psi_i\rangle\langle\psi_i|]
= \frac{1}{i\hbar}[H,\rho].
\]
因此无论是纯态还是混合态，密度算符均满足同一运动方程。这就是冯诺依曼方程（量子刘维尔方程）。
\end{solution}

\begin{question}
对沿 $y$ 方向自旋向下的电子重复例题 12.1 的计算。
\end{question}
\begin{solution}
取 $S_z$ 表象，泡利矩阵为
\[
\sigma_x = \begin{pmatrix}0&1\\1&0\end{pmatrix},\quad
\sigma_y = \begin{pmatrix}0&-i\\i&0\end{pmatrix},\quad
\sigma_z = \begin{pmatrix}1&0\\0&-1\end{pmatrix}.
\]
沿 $y$ 方向自旋向下的本征态满足 $\mathbf S\cdot\hat y\,\chi_{-y} = -\frac{\hbar}{2}\chi_{-y}$，其中 $\mathbf S = \frac{\hbar}{2}\boldsymbol\sigma$，即
\[
(\hat y\cdot\boldsymbol\sigma)\chi_{-y} = -\chi_{-y},\qquad
\hat y\cdot\boldsymbol\sigma = \sigma_y = \begin{pmatrix}0&-i\\i&0\end{pmatrix}.
\]
解本征方程：
\[
\begin{pmatrix}0&-i\\i&0\end{pmatrix}\begin{pmatrix}a\\b\end{pmatrix} = -\begin{pmatrix}a\\b\end{pmatrix}
\;\Longrightarrow\;
\begin{cases}
-ib = -a,\\
ia = -b,
\end{cases}
\]
两式等价，给出 $a = ib$。归一化得 $|a|^2+|b|^2 = 2|b|^2 = 1$，故取
\[
\chi_{-y} = \frac{1}{\sqrt{2}}\begin{pmatrix}1\\-i\end{pmatrix}.
\]

密度矩阵为
\[
\rho = \chi_{-y}\chi_{-y}^\dagger
= \frac{1}{2}\begin{pmatrix}1\\-i\end{pmatrix}\begin{pmatrix}1&i\end{pmatrix}
= \frac{1}{2}\begin{pmatrix}1 & i\\ -i & 1\end{pmatrix}.
\]

各方向自旋期望值用 $\langle S_i\rangle = \frac{\hbar}{2}\operatorname{Tr}(\rho\sigma_i)$ 计算：

对于 $S_x$：
\[
\rho\sigma_x = \frac{1}{2}\begin{pmatrix}1&i\\-i&1\end{pmatrix}\begin{pmatrix}0&1\\1&0\end{pmatrix}
= \frac{1}{2}\begin{pmatrix}i&1\\1&-i\end{pmatrix},
\]
\[
\operatorname{Tr}(\rho\sigma_x) = \frac{1}{2}(i - i) = 0
\;\Longrightarrow\; \langle S_x\rangle = 0.
\]

对于 $S_y$：
\[
\rho\sigma_y = \frac{1}{2}\begin{pmatrix}1&i\\-i&1\end{pmatrix}\begin{pmatrix}0&-i\\i&0\end{pmatrix}
= \frac{1}{2}\begin{pmatrix}-1&0\\0&-1\end{pmatrix} = -\frac{1}{2}I,
\]
\[
\operatorname{Tr}(\rho\sigma_y) = -1
\;\Longrightarrow\; \langle S_y\rangle = \frac{\hbar}{2}(-1) = -\frac{\hbar}{2}.
\]

对于 $S_z$：
\[
\rho\sigma_z = \frac{1}{2}\begin{pmatrix}1&i\\-i&1\end{pmatrix}\begin{pmatrix}1&0\\0&-1\end{pmatrix}
= \frac{1}{2}\begin{pmatrix}1&-i\\-i&-1\end{pmatrix},
\]
\[
\operatorname{Tr}(\rho\sigma_z) = \frac{1}{2}(1-1) = 0
\;\Longrightarrow\; \langle S_z\rangle = 0.
\]

综上，
\[
\boxed{\langle S_x\rangle = 0,\qquad \langle S_y\rangle = -\frac{\hbar}{2},\qquad \langle S_z\rangle = 0.}
\]

沿任意方向 $\hat n = (\sin\theta\cos\phi,\;\sin\theta\sin\phi,\;\cos\theta)$ 测得 $+\hbar/2$ 的概率为
\[
P_+(\hat n) = \operatorname{Tr}(\rho P_{+\hat n}),
\]
其中 $P_{+\hat n} = \frac{1}{2}(I+\hat n\cdot\boldsymbol\sigma)$ 是沿 $\hat n$ 方向自旋向上的投影算符。计算：
\[
\rho P_{+\hat n} = \frac{1}{4}\bigl(I-i\sigma_y\bigr)\bigl(I+\hat n\cdot\boldsymbol\sigma\bigr),
\]
其中利用了 $\rho = \frac{1}{2}(I-i\sigma_y)$（因 $\chi_{-y}$ 对应的 Bloch 矢量为 $-\hat y$）。展开并取迹，利用 $\operatorname{Tr}\sigma_i=0$，$\operatorname{Tr}(\sigma_i\sigma_j)=2\delta_{ij}$：
\[
\begin{aligned}
\operatorname{Tr}(\rho P_{+\hat n})
&= \frac{1}{4}\operatorname{Tr}\bigl[I - i\sigma_y + \hat n\cdot\boldsymbol\sigma - i(\hat n\cdot\boldsymbol\sigma)\sigma_y\bigr]\\[2mm]
&= \frac{1}{4}\bigl(2 + 0 + 0 - i\cdot 2 n_y\bigr)
= \frac{1}{2}(1 - n_y).
\end{aligned}
\]
由于 $\hat n$ 的 $y$ 分量为 $n_y = \sin\theta\sin\phi$，且 $\hat y\cdot\hat n = n_y$，故
\[
\boxed{P_+(\hat n) = \frac{1}{2}\bigl(1 - \hat y\cdot\hat n\bigr).}
\]
特别地，当 $\hat n = \hat y$ 时 $P_+ = 0$（测到向上的概率为零，与预期一致）；当 $\hat n = -\hat y$ 时 $P_+ = 1$（测到向下的概率为 1）。
\end{solution}

\begin{question}
\begin{enumerate}
\item 证明式 (12.31)--式 (12.34)。
\item 证明 $\operatorname{Tr}(\rho^2)\le 1$，且仅当 $\rho$ 表示纯态时才等于 1。
\item 证明：当且仅当 $\rho$ 表示纯态时，$\rho^2=\rho$。
\end{enumerate}

\par\smallskip
\noindent{\kaishu 为避免查阅正文，本题中引用的公式为：式 (12.31--12.34)：密度矩阵的判据：厄米、迹为 $1$、本征值非负；任意观测量平均值仍为 $\langle Q\rangle=\operatorname{Tr}(\rho Q)$。}

\end{question}
\begin{solution}
(a) 密度算符的性质（式 (12.31)--(12.34)）：

\textbf{厄米性（式 (12.31)）。} 由定义 $\rho = \sum_i p_i |\psi_i\rangle\langle\psi_i|$，取共轭得 $\rho^\dagger = \rho$，故密度算符是厄米算符。厄米算符的本征值为实数，可在某正交基下对角化。

\textbf{迹为 1（式 (12.32)）。}
\[
\operatorname{Tr}\rho = \sum_i p_i \operatorname{Tr}(|\psi_i\rangle\langle\psi_i|) = \sum_i p_i = 1.
\]

\textbf{本征值非负（式 (12.33)）。} 在 $\rho$ 的对角化基中，$\rho = \sum_n \lambda_n |n\rangle\langle n|$，其中 $\lambda_n$ 是本征值。由于 $\rho$ 是概率系综的表示：对任意态 $|\phi\rangle$，
\[
\langle\phi|\rho|\phi\rangle = \sum_i p_i |\langle\phi|\psi_i\rangle|^2 \ge 0,
\]
因此 $\rho$ 是正定算符，所有本征值 $\lambda_n \ge 0$。

\textbf{期望值公式（式 (12.34)）。} 任意可观测量 $Q$ 的期望值为
\[
\langle Q\rangle = \sum_i p_i \langle\psi_i|Q|\psi_i\rangle
= \sum_i p_i \operatorname{Tr}(Q|\psi_i\rangle\langle\psi_i|)
= \operatorname{Tr}\Bigl(Q\sum_i p_i|\psi_i\rangle\langle\psi_i|\Bigr)
= \operatorname{Tr}(\rho Q).
\]

(b) 在 $\rho$ 的对角化基中，设本征值为 $\lambda_n \ge 0$，且 $\sum_n \lambda_n = 1$。则
\[
\operatorname{Tr}(\rho^2) = \sum_n \lambda_n^2.
\]
利用柯西-施瓦茨不等式或简单的代数关系：
\[
\Bigl(\sum_n \lambda_n\Bigr)^2 = \sum_n \lambda_n^2 + 2\sum_{m<n} \lambda_m\lambda_n \ge \sum_n \lambda_n^2,
\]
等号成立当且仅当至多一个 $\lambda_n$ 非零。由于 $\sum_n \lambda_n = 1$，故
\[
\sum_n \lambda_n^2 \le \Bigl(\sum_n \lambda_n\Bigr)^2 = 1,
\]
即
\[
\boxed{\operatorname{Tr}(\rho^2) \le 1.}
\]

等号成立的条件：$\operatorname{Tr}(\rho^2)=1$ 意味着 $\sum_n \lambda_n^2 = 1 = (\sum_n \lambda_n)^2$，这只有当一个 $\lambda_k = 1$ 而其余 $\lambda_{n\neq k} = 0$ 时才成立。此时 $\rho = |k\rangle\langle k|$ 是投影到单维子空间的算符，即纯态。

(c) 分两个方向证明：

\textbf{充分性：} 若 $\rho$ 是纯态，即 $\rho = |\psi\rangle\langle\psi|$，则
\[
\rho^2 = |\psi\rangle\langle\psi|\psi\rangle\langle\psi| = |\psi\rangle\langle\psi| = \rho,
\]
其中用到 $\langle\psi|\psi\rangle=1$。故 $\rho^2 = \rho$。

\textbf{必要性：} 若 $\rho^2 = \rho$，则 $\rho$ 是投影算符。在对角化基中 $\rho = \sum_n \lambda_n |n\rangle\langle n|$，则
\[
\rho^2 = \sum_n \lambda_n^2 |n\rangle\langle n| = \sum_n \lambda_n |n\rangle\langle n| = \rho,
\]
因此每个本征值满足 $\lambda_n^2 = \lambda_n$，即 $\lambda_n = 0$ 或 $1$。又因为 $\operatorname{Tr}\rho = \sum_n \lambda_n = 1$，所以恰好有一个本征值为 1，其余全为 0。故 $\rho = |k\rangle\langle k|$ 是纯态。

综上，
\[
\boxed{\rho^2 = \rho \;\Longleftrightarrow\; \rho \text{ 是纯态}.}
\]

作为补充，$\operatorname{Tr}(\rho^2)$ 称为纯度（purity），是衡量系统混合程度的常用量：
\[
\gamma \equiv \operatorname{Tr}(\rho^2),\qquad \frac{1}{d} \le \gamma \le 1,
\]
其中 $d$ 为 Hilbert 空间维数。$\gamma=1$ 对应纯态，$\gamma=1/d$ 对应最大混合态。
\end{solution}

\begin{question}
\begin{enumerate}
\item 构造电子状态处于沿 $x$ 轴自旋向上（几率为 $1/3$）或沿 $y$ 轴自旋向下（几率为 $2/3$）的密度矩阵。
\item 对 (a) 中的电子，计算 $\langle S_y\rangle$。
\end{enumerate}
\end{question}
\begin{solution}
(a) 首先构造每个纯态的密度矩阵。在 $S_z$ 表象中：

沿 $x$ 轴自旋向上的本征态为
\[
\chi_{+x} = \frac{1}{\sqrt{2}}\begin{pmatrix}1\\1\end{pmatrix},
\]
相应的投影算符为
\[
\rho_{+x} = \chi_{+x}\chi_{+x}^\dagger
= \frac{1}{2}\begin{pmatrix}1\\1\end{pmatrix}\begin{pmatrix}1&1\end{pmatrix}
= \frac{1}{2}\begin{pmatrix}1&1\\1&1\end{pmatrix}.
\]
利用泡利矩阵可写为 $\rho_{+x} = \frac{1}{2}(I+\sigma_x)$。

沿 $y$ 轴自旋向下的本征态为
\[
\chi_{-y} = \frac{1}{\sqrt{2}}\begin{pmatrix}1\\-i\end{pmatrix},
\]
相应的投影算符为
\[
\rho_{-y} = \chi_{-y}\chi_{-y}^\dagger
= \frac{1}{2}\begin{pmatrix}1\\-i\end{pmatrix}\begin{pmatrix}1&i\end{pmatrix}
= \frac{1}{2}\begin{pmatrix}1&i\\-i&1\end{pmatrix}.
\]
利用泡利矩阵可写为 $\rho_{-y} = \frac{1}{2}(I-\sigma_y)$（因为 $-\hat y$ 方向的 Bloch 矢量为 $-\sigma_y$）。

混合态的密度算符为两个纯态投影的概率加权和：
\[
\rho = \frac{1}{3}\rho_{+x} + \frac{2}{3}\rho_{-y}
= \frac{1}{3}\cdot\frac{1}{2}(I+\sigma_x) + \frac{2}{3}\cdot\frac{1}{2}(I-\sigma_y)
= \frac{1}{2}\left(I + \frac{1}{3}\sigma_x - \frac{2}{3}\sigma_y\right).
\]

代入泡利矩阵的显式形式：
\[
I = \begin{pmatrix}1&0\\0&1\end{pmatrix},\quad
\sigma_x = \begin{pmatrix}0&1\\1&0\end{pmatrix},\quad
\sigma_y = \begin{pmatrix}0&-i\\i&0\end{pmatrix},
\]
得
\[
\rho = \frac{1}{2}\left[
\begin{pmatrix}1&0\\0&1\end{pmatrix}
+ \frac{1}{3}\begin{pmatrix}0&1\\1&0\end{pmatrix}
- \frac{2}{3}\begin{pmatrix}0&-i\\i&0\end{pmatrix}
\right].
\]
计算各矩阵元：
\begin{itemize}
\item $\rho_{11} = \frac{1}{2}(1 + 0 + 0) = \frac{1}{2}$
\item $\rho_{12} = \frac{1}{2}\left(0 + \frac{1}{3}\cdot 1 - \frac{2}{3}\cdot(-i)\right) = \frac{1}{2}\left(\frac{1}{3} + \frac{2i}{3}\right) = \frac{1}{6} + \frac{i}{3}$
\item $\rho_{21} = \frac{1}{2}\left(0 + \frac{1}{3}\cdot 1 - \frac{2}{3}\cdot i\right) = \frac{1}{2}\left(\frac{1}{3} - \frac{2i}{3}\right) = \frac{1}{6} - \frac{i}{3}$
\item $\rho_{22} = \frac{1}{2}(1 + 0 + 0) = \frac{1}{2}$
\end{itemize}
因此
\[
\boxed{\rho = \frac{1}{2}\begin{pmatrix}
1 & \dfrac{1}{3} + \dfrac{2i}{3}\\[6pt]
\dfrac{1}{3} - \dfrac{2i}{3} & 1
\end{pmatrix}.}
\]

(b) 期望值 $\langle S_y\rangle$ 可通过多种方法计算。

\textbf{方法一：} 利用 $\langle S_y\rangle = \operatorname{Tr}(\rho S_y)$，其中 $S_y = \frac{\hbar}{2}\sigma_y$：
\[
\langle S_y\rangle = \frac{\hbar}{2}\operatorname{Tr}(\rho\sigma_y).
\]
由 $\rho$ 的 Bloch 矢量形式 $\rho = \frac{1}{2}(I + \mathbf a\cdot\boldsymbol\sigma)$，其中 $\mathbf a = \left(\frac{1}{3}, -\frac{2}{3}, 0\right)$。则
\[
\rho\sigma_y = \frac{1}{2}(I + \mathbf a\cdot\boldsymbol\sigma)\sigma_y,
\]
利用 $\operatorname{Tr}\sigma_y = 0$ 和 $\operatorname{Tr}(\sigma_i\sigma_j)=2\delta_{ij}$：
\[
\operatorname{Tr}(\rho\sigma_y) = \frac{1}{2}\operatorname{Tr}(\sigma_y) + \frac{1}{2}\sum_j a_j\operatorname{Tr}(\sigma_j\sigma_y) = \frac{1}{2}\cdot 0 + \frac{1}{2}\cdot 2a_y = a_y = -\frac{2}{3}.
\]
因此
\[
\boxed{\langle S_y\rangle = \frac{\hbar}{2}\left(-\frac{2}{3}\right) = -\frac{\hbar}{3}.}
\]

\textbf{方法二（直接矩阵乘法验证）：}
\[
\rho\sigma_y = \frac{1}{2}\begin{pmatrix}
1 & \frac{1+2i}{3} \\[4pt]
\frac{1-2i}{3} & 1
\end{pmatrix}
\begin{pmatrix}
0 & -i \\ i & 0
\end{pmatrix}
= \frac{1}{2}\begin{pmatrix}
\frac{i(1+2i)}{3} & -i \\[4pt]
i & -\frac{i(1-2i)}{3}
\end{pmatrix},
\]
取迹：
\[
\operatorname{Tr}(\rho\sigma_y) = \frac{1}{2}\left(\frac{i(1+2i)}{3} - \frac{i(1-2i)}{3}\right)
= \frac{1}{2}\cdot\frac{i(1+2i-1+2i)}{3}
= \frac{1}{2}\cdot\frac{4i^2}{3} = -\frac{2}{3},
\]
与上述结果一致。
\end{solution}

\begin{question}
\begin{enumerate}
\item 证明自旋 $1/2$ 粒子最常见的密度矩阵可以用 3 个实数 $(a_1,a_2,a_3)$ 表示为
\[
\rho=\frac12\begin{pmatrix}1+a_3&a_1-ia_2\\a_1+ia_2&1-a_3\end{pmatrix}=\frac12(1+\mathbf a\cdot\boldsymbol\sigma),
\]
其中 $\sigma_1,\sigma_2,\sigma_3$ 是泡利矩阵。提示：它必须是厄米的，且迹必须为 1。
\item 证明当且仅当 $|\mathbf a|=1$ 时 $\rho$ 表示纯态；当 $|\mathbf a|<1$ 时表示混合态。
\item 若布洛赫矢量的端点分别位于北极 $(0,0,1)$、球心 $(0,0,0)$ 和南极 $(0,0,-1)$，对 $S_z$ 测量得到 $+\hbar/2$ 的几率是多少？
\item 若纯态的布洛赫矢量位于赤道上且方位角为 $\phi$，求系统的旋量 $\chi$。
\end{enumerate}
\end{question}
\begin{solution}
(a) 任意 $2\times2$ 厄米矩阵 $M$ 可以展开为四个实参数：
\[
M = c_0 I + c_1\sigma_1 + c_2\sigma_2 + c_3\sigma_3,
\]
其中 $c_0,c_1,c_2,c_3$ 为实数，这是因为 $I,\sigma_1,\sigma_2,\sigma_3$ 构成 $2\times2$ 厄米矩阵空间的一组完备基（四个基矩阵线性独立且均为厄米矩阵）。

密度算符 $\rho$ 必须是厄米的，故可写为上述形式。迹为 1 的条件给出：
\[
\operatorname{Tr}\rho = \operatorname{Tr}(c_0 I) + \sum_{j=1}^3 c_j\operatorname{Tr}\sigma_j = 2c_0 = 1 \;\Longrightarrow\; c_0 = \frac12.
\]
记 $\mathbf a = (a_1,a_2,a_3) = 2(c_1,c_2,c_3)$，则
\[
\rho = \frac12 I + \frac12\sum_{j=1}^3 a_j\sigma_j = \frac12(I + \mathbf a\cdot\boldsymbol\sigma).
\]

在 $S_z$ 表象中写出矩阵形式。泡利矩阵为
\[
\sigma_1 = \begin{pmatrix}0&1\\1&0\end{pmatrix},\quad
\sigma_2 = \begin{pmatrix}0&-i\\i&0\end{pmatrix},\quad
\sigma_3 = \begin{pmatrix}1&0\\0&-1\end{pmatrix}.
\]
代入得
\[
\rho = \frac12\left[
\begin{pmatrix}1&0\\0&1\end{pmatrix}
+ a_1\begin{pmatrix}0&1\\1&0\end{pmatrix}
+ a_2\begin{pmatrix}0&-i\\i&0\end{pmatrix}
+ a_3\begin{pmatrix}1&0\\0&-1\end{pmatrix}
\right].
\]
计算各矩阵元：
\begin{itemize}
\item $\rho_{11} = \frac12(1 + a_3)$
\item $\rho_{12} = \frac12(a_1 - i a_2)$
\item $\rho_{21} = \frac12(a_1 + i a_2)$
\item $\rho_{22} = \frac12(1 - a_3)$
\end{itemize}
因此
\[
\boxed{\rho = \frac12\begin{pmatrix}
1+a_3 & a_1-ia_2\\
a_1+ia_2 & 1-a_3
\end{pmatrix}
= \frac12(I+\mathbf a\cdot\boldsymbol\sigma).}
\]

(b) 求 $\rho$ 的本征值。设 $\mathbf a = |\mathbf a|\hat n$，则 $\hat n\cdot\boldsymbol\sigma$ 的本征值为 $\pm 1$。因此
\[
\rho = \frac12(I + |\mathbf a|\,\hat n\cdot\boldsymbol\sigma)
\]
的本征值为
\[
\lambda_\pm = \frac12(1 \pm |\mathbf a|).
\]
密度算符的本征值必须非负：$\lambda_\pm \ge 0 \;\Longrightarrow\; |\mathbf a| \le 1$。

纯态条件 $\rho^2 = \rho$ 等价于 $\lambda_\pm^2 = \lambda_\pm$。代入本征值：
\[
\frac14(1\pm|\mathbf a|)^2 = \frac12(1\pm|\mathbf a|).
\]
两边乘以 4：
\[
(1\pm|\mathbf a|)^2 = 2(1\pm|\mathbf a|).
\]
令 $x = 1\pm|\mathbf a|$，则 $x^2 = 2x$，解得 $x=0$ 或 $x=2$。对应地：
\begin{itemize}
\item $1+|\mathbf a| = 2$ \;\Longrightarrow\; $|\mathbf a| = 1$，此时 $\lambda_+ = 1$, $\lambda_- = 0$。
\item $1-|\mathbf a| = 0$ \;\Longrightarrow\; $|\mathbf a| = 1$，此时 $\lambda_+ = 1$, $\lambda_- = 0$（与上同）。
\end{itemize}
因此当且仅当 $|\mathbf a| = 1$ 时 $\rho$ 有一个本征值为 1、另一个为 0，即 $\rho$ 是纯态。

当 $|\mathbf a| < 1$ 时，两个本征值均为正且小于 1，$\operatorname{Tr}(\rho^2) = \lambda_+^2 + \lambda_-^2 = \frac12(1+|\mathbf a|^2) < 1$，表示混合态。

综上，
\[
\boxed{|\mathbf a| = 1 \;\Longleftrightarrow\; \text{纯态},\qquad
|\mathbf a| < 1 \;\Longleftrightarrow\; \text{混合态}.}
\]

(c) 对于 $S_z$ 测量，得到 $+\hbar/2$ 的概率是 $|\langle\uparrow|\rho|\uparrow\rangle|$（实际上是 $\rho$ 的 $(1,1)$ 矩阵元）。由 (a) 的表达式：
\[
P_+ = \rho_{11} = \frac{1+a_3}{2}.
\]
\begin{itemize}
\item 北极 $\mathbf a = (0,0,1)$：$a_3 = 1$，$P_+ = \frac{1+1}{2} = 1$（自旋完全沿 $+z$ 方向）。
\item 球心 $\mathbf a = (0,0,0)$：$a_3 = 0$，$P_+ = \frac{1}{2}$（完全非极化，等概率混合）。
\item 南极 $\mathbf a = (0,0,-1)$：$a_3 = -1$，$P_+ = \frac{1-1}{2} = 0$（自旋完全沿 $-z$ 方向）。
\end{itemize}
故
\[
\boxed{P_+(\text{北极})=\;1,\qquad P_+(\text{球心})=\frac12,\qquad P_+(\text{南极})=\;0.}
\]

(d) 纯态 $|\mathbf a|=1$ 且位于赤道上，即 $\mathbf a = (\cos\phi,\;\sin\phi,\;0)$。对应的密度矩阵为
\[
\rho = \frac12\begin{pmatrix}
1 & \cos\phi - i\sin\phi\\
\cos\phi + i\sin\phi & 1
\end{pmatrix}
= \frac12\begin{pmatrix}
1 & e^{-i\phi}\\
e^{i\phi} & 1
\end{pmatrix}.
\]
由于是纯态，$\rho = \chi\chi^\dagger$。设旋量为
\[
\chi = \begin{pmatrix}a\\b\end{pmatrix},\quad |a|^2+|b|^2=1,
\]
则
\[
\chi\chi^\dagger = \begin{pmatrix}|a|^2 & ab^*\\ a^*b & |b|^2\end{pmatrix}
= \frac12\begin{pmatrix}1 & e^{-i\phi}\\ e^{i\phi} & 1\end{pmatrix}.
\]
比较得 $|a|^2 = |b|^2 = \frac12$，且 $ab^* = \frac12 e^{-i\phi}$。可取 $a = \frac{1}{\sqrt{2}}$，$b = \frac{1}{\sqrt{2}}e^{i\phi}$（整体相位自由），因此
\[
\boxed{\chi = \frac{1}{\sqrt{2}}\begin{pmatrix}1\\ e^{i\phi}\end{pmatrix},}
\]
整体相位可任意选择。这一旋量对应沿赤道方向 $\hat n = (\cos\phi,\sin\phi,0)$ 的自旋向上态。
\end{solution}
